% TOP_PROBLEM: {"id": 2965, "title": "Kirby Problem 4.89", "category_id": 7, "category": {"id": 7, "name": "topology", "display_name": "Topology", "description": "Properties preserved under continuous deformations.", "slug": "topology", "order_index": 7, "created_at": "2026-07-31T15:26:25.671Z"}, "status": "open", "problem_number": "KP-4.89", "set_id": 11, "statusQualification": "Compact nonempty-boundary setting follows the source bounded-case discussion; relative and unmarked uniqueness conventions are made explicit."}
% FIRST_POSTED: 2026-10-01
% ENTRY_KIND: statement_only
% UnsolvedMath ID 2965; source catalogue code KP-4.89
% !TeX program = lualatex
% Definition and sources only; this file is not an LLM solution attempt.
\documentclass[11pt,a4paper]{article}
\usepackage[margin=25mm]{geometry}
\usepackage{fontspec}
\setmainfont{lmroman10-regular.otf}[BoldFont=lmroman10-bold.otf,
 ItalicFont=lmroman10-italic.otf,BoldItalicFont=lmroman10-bolditalic.otf]
\usepackage{amsmath,amssymb,mathtools}
\usepackage{unicode-math}
\setmathfont{latinmodern-math.otf}
\AtBeginDocument{\let\setminus\smallsetminus}
\usepackage{xurl,hyperref}
\hypersetup{colorlinks=true,urlcolor=blue}
\setcounter{secnumdepth}{0}
\setlength{\parindent}{0pt}
\setlength{\parskip}{5pt}
\setlength{\emergencystretch}{3em}
\begin{document}
\section*{Kirby Problem 4.89}\label{2965}
{\small UnsolvedMath ID 2965 · KP-4.89 · Topology · Kirby's Problems in Low-Dimensional Topology}
\subsection{Definitions and mathematical statement}
A smooth circle action on a smooth $4$-manifold $X$ is a smooth
map $a:S^1\times X\to X$ with $a(1,x)=x$ and
$a(uv,x)=a(u,a(v,x))$. It is effective if the only circle element
fixing all of $X$ is $1$. An action preserves the boundary.
Consider simply connected compact smooth $4$-manifolds with nonempty
boundary, in the bounded setting of the source.

Part (a) asks which diffeomorphism types of such manifolds admit
an effective circle action. Part (b) asks, for each eligible $X$,
to classify its actions up to smooth conjugacy. Conjugacy means
that a diffeomorphism $F:X\to X'$ satisfies
$F(a(u,x))=a'(u,F(x))$ for every $u\in S^1$ and $x\in X$;
the circle parameter is kept fixed.

Part (c) starts with a prescribed smooth circle action on
$\partial X$ and asks for necessary and sufficient conditions for
an extension to a smooth action on $X$, and for uniqueness of
that extension. For a fixed boundary marking, uniqueness is
understood up to equivariant diffeomorphism restricting to the
identity on the boundary; forgetting the marking gives the
corresponding conjugacy question.

The extension must satisfy the circle group law throughout the
interior. Extending each boundary diffeomorphism separately, or
extending only the loop it represents in a diffeomorphism group,
does not necessarily supply a circle action. Fixed points and
exceptional orbits in the interior are allowed.
\subsection{Short English statement}
Classify effective smooth circle actions on simply connected compact four-manifolds with boundary, including existence, conjugacy and extension of prescribed boundary actions.
\subsection{Sources}
\begin{itemize}
\item[S1] ULAM.AI and the UnsolvedMath Contributors, supplied problems.json, record 2965 (KP-4.89), Kirby's Problems in Low-Dimensional Topology. Catalogue identity and source transcription; CC BY 4.0.
\url{https://huggingface.co/datasets/ulamai/UnsolvedMath}.
\item[S2] K3: A New Problem List in Low-Dimensional Topology, Problem 4.89. Expanded formulation of the supplied catalogue statement.
\url{https://math.berkeley.edu/sites/default/files/surv-295-ruberman-watermarked-author-pdf.pdf}.
\end{itemize}
\end{document}
